% Zone „Das kennst du schon“ – aus bank/bruchrechnung/zone.jsonl (zusammenbau v0.1)
\einheitenkopf[zone][Kennst du schon]{Das kennst du schon}
\setcounter{aufgabe}{0}
%% TODO Grundvorstellungs-Aufgabe als erste Hauptnummer der Zone (2.8) fehlt in der Bank

%% TODO Titel: Ich-kann-Satz fehlt in der Bank (Platzhalter: Kettenname) – Nr. 1
%% TODO Anweisung: ein Satz über der ersten Teilaufgabe fehlt in der Bank – Nr. 1
\begin{aufgabe}{Bruch als Anteil lesen, am Streifen einzeichnen (drei Achtel markieren)}
\swa{Welcher Anteil ist gefärbt? \leerfeld}{\bruchrechteck{2}{5}}
\swa{Der Streifen hat 12 Teile. Färbe $\frac{2}{3}$ davon.}{\bruchrechteck{0}{12}}
\end{aufgabe}

%% TODO Titel: Ich-kann-Satz fehlt in der Bank (Platzhalter: Kettenname) – Nr. 2
%% TODO Anweisung: ein Satz über der ersten Teilaufgabe fehlt in der Bank – Nr. 2
\begin{aufgabe}{Kürzen und Erweitern; gleichwertige Brüche (zwei Drittel gleich vier Sechstel)}
%% TODO Darstellung für schwach fehlt in der Bank (bruchrechnung-zone-f2-v1; Abschnitt „Für schwache Schüler“)
\swz{Erweitere $\frac{2}{5}$ mit 3.}{}
%% TODO Darstellung für schwach fehlt in der Bank (bruchrechnung-zone-f2-v3; Abschnitt „Für schwache Schüler“)
\swz{Kürze $\frac{24}{36}$ so weit wie möglich.}{}
\end{aufgabe}

%% TODO Titel: Ich-kann-Satz fehlt in der Bank (Platzhalter: Kettenname) – Nr. 3
%% TODO Anweisung: ein Satz über der ersten Teilaufgabe fehlt in der Bank – Nr. 3
\begin{aufgabe}{Vielfache und Teiler (kgV von vier und sechs, ggT von zwölf und achtzehn)}
%% TODO Darstellung für schwach fehlt in der Bank (bruchrechnung-zone-f3-v1; Abschnitt „Für schwache Schüler“)
\swz{Nenne die ersten vier Vielfachen von 3.}{}
%% TODO Darstellung für schwach fehlt in der Bank (bruchrechnung-zone-f3-v3; Abschnitt „Für schwache Schüler“)
\swz{Größter gemeinsamer Teiler von 16 und 24?}{}
\end{aufgabe}

%% TODO Titel: Ich-kann-Satz fehlt in der Bank (Platzhalter: Kettenname) – Nr. 4
%% TODO Anweisung: ein Satz über der ersten Teilaufgabe fehlt in der Bank – Nr. 4
\begin{aufgabe}{Stellenwerte der Dezimalzahlen (Zehntel, Hundertstel; 0,7 = 0,70)}
%% TODO Darstellung für schwach fehlt in der Bank (bruchrechnung-zone-f4-v1; Abschnitt „Für schwache Schüler“)
\swz{4 Zehntel als Dezimalzahl}{}
%% TODO Darstellung für schwach fehlt in der Bank (bruchrechnung-zone-f4-v3; Abschnitt „Für schwache Schüler“)
\swz{2 Einer, 5 Zehntel und 8 Hundertstel als Dezimalzahl}{}
\end{aufgabe}

%% TODO Titel: Ich-kann-Satz fehlt in der Bank (Platzhalter: Kettenname) – Nr. 5
%% TODO Anweisung: ein Satz über der ersten Teilaufgabe fehlt in der Bank – Nr. 5
\begin{aufgabe}{Schriftliches Rechnen mit natürlichen Zahlen, Einmaleins}
%% TODO Darstellung für schwach fehlt in der Bank (bruchrechnung-zone-f5-v1; Abschnitt „Für schwache Schüler“)
\swz{$7 \cdot 8$}{}
%% TODO Darstellung für schwach fehlt in der Bank (bruchrechnung-zone-f5-v3; Abschnitt „Für schwache Schüler“)
\swz{Schriftlich: $3\,475\text{ g} + 1\,268\text{ g}$}{}
\end{aufgabe}

%% TODO Titel: Ich-kann-Satz fehlt in der Bank (Platzhalter: Kettenname) – Nr. 6
%% TODO Anweisung: ein Satz über der ersten Teilaufgabe fehlt in der Bank – Nr. 6
\begin{aufgabe}{Bruch ↔ Dezimalzahl bei einfachen Brüchen (ein Halb, drei Viertel)}
%% TODO Darstellung für schwach fehlt in der Bank (bruchrechnung-zone-f6-v1; Abschnitt „Für schwache Schüler“)
\swz{$\frac{3}{10}$ als Dezimalzahl}{}
%% TODO Darstellung für schwach fehlt in der Bank (bruchrechnung-zone-f6-v3; Abschnitt „Für schwache Schüler“)
\swz{$\frac{3}{5}$ als Dezimalzahl}{}
\end{aufgabe}

%% TODO Titel: Ich-kann-Satz fehlt in der Bank (Platzhalter: Kettenname) – Nr. 7
%% TODO Anweisung: ein Satz über der ersten Teilaufgabe fehlt in der Bank – Nr. 7
\begin{aufgabe}{Kürzen und Erweitern; gleichwertige Brüche (zwei Drittel gleich vier Sechstel) – Fehler finden}
\swz[3]{Lea erweitert: $\frac{3}{5} = \frac{3}{10}$. Finde den Fehler.}{}
\end{aufgabe}

%% TODO Titel: Ich-kann-Satz fehlt in der Bank (Platzhalter: Kettenname) – Nr. 8
%% TODO Anweisung: ein Satz über der ersten Teilaufgabe fehlt in der Bank – Nr. 8
\begin{aufgabe}{Kürzen und Erweitern; gleichwertige Brüche (zwei Drittel gleich vier Sechstel) – selbst rechnen}
%% TODO Darstellung für schwach fehlt in der Bank (bruchrechnung-zone-f2-v6; Abschnitt „Für schwache Schüler“)
\swz{Schreibe $\frac{4}{9}$ mit dem Nenner 27.}{}
\end{aufgabe}
